\scanpage{421}
\mspage{160}
$\cdot$ Je vais prouver $1,2,1',2'$.

1) Considérons le diagramme
\diagram{\begin{tikzcd}[ampersand replacement=\&,column sep=small]
X=(X_{\unsure{\underline{\mathrm{Tz}}|\underline{\mathrm{Ty}}|\underline{\mathrm{Iz}}|\underline{\mathrm{Iy}}}})_e\arrow[r,hook,"i"]\arrow[dr,dashed] \&
X_{\underline{\mathrm{Tz}}}=X\times\underline{\mathrm{Tz}}\arrow[rr,"j"]\arrow[dr] \&\&\underline{\mathrm{Ch}}(X)\arrow[dl]\\
\& e\arrow[r,dashed] \&\underline{\mathrm{Tz}}\&{,}
\end{tikzcd}}
où $X_{\underline{\mathrm{Tz}}}$ et $\underline{\mathrm{Ch}}(X)$ sont considérées comme catégories Cat-fibrées sur $\underline{\mathrm{Tz}}$, et le foncteur $j$ est déduit des foncteurs canoniques
\[
X\longrightarrow\underline{\mathrm{Ch}}(X)_\tau\overset{\mathrm{def}}{=}\underline{\mathrm{Ch}}_\tau(X),
\]
associant à \deleted{chaque} $x\in\mathrm{Ob}\,X$ le chemin \deleted{vide} $\unsure{p|\varphi}_\tau^*(x)$, où $p_\tau:\tau\mathrel{\unsure{\text{---}|\to}}e$ est la projection de $\tau$ vers l'élément unité (et élément final) $e=1_{\underline{\mathrm{Tz}}}$ de $\underline{\mathrm{Tz}}$, et $\underline{\mathrm{Ch}}_e(X)$ est identifié à $X$. Il faut vérifier le

\uline{Lemme}\quad Pour tout $\tau\in\mathrm{Tz}$, le foncteur précédent est une Hot-équivalence (et même une \deleted{\unsure{équivale}} homotopisme).

Mais on sait que pour deux catégories $\unsure{\mathscr X|\mathcal X},\mathcal C$ lisses sur une \unsure{autre|même} $\unsure{T_0|T}$, tout $T$-morphisme de $\mathscr X$ dans $\mathcal C$ qui induit des Hot-équivalences sur les fibres, est une Hot-équivalence. Ainsi $j$ est
% archive sheet number: 420
